
\color{unidarkgreen}{\subsection{Пусковой орган устройства резервирования при отказе выключателя при повреждениях на стороне НН}}
\color{black}

\begin{enumerate}[label=\arabic{section}.\arabic{subsection}.\arabic*,  align=left, labelsep=4pt, leftmargin=0pt, itemindent=57pt]

\item
Пусковой орган устройства резервирования при отказе выключателя при повреждениях на стороне НН (ПО УРОВ НН) предназначен для пуска УРОВ выключателей сторон ВН и СН автотрансформатора при недостаточной чувствительности своих измерительных органов по току при повреждениях на стороне НН автотрансформатора.

\item
Логика работы ПО УРОВ НН выполнена с контролем протекания тока по стороне НН автотрансформатора. Значение по току срабатывания задается параметром \textbf{«Iср УРОВ В»}. Функция УРОВ действует с регулируемой задержкой по времени \textbf{«Тср повт. УРОВ»} на повторное отключение своего выключателя, с задержкой \textbf{«Тср УРОВ»} на отключение смежных выключателей.

\item
Предусмотрен подхват по току, включение которого производится программной накладкой \textbf{«ПО УРОВ подхв.»}.

\item
Алгоритм работы ПО УРОВ НН приведен на рисунке \ref{pic:UROVNN}. Уставки ПО УРОВ НН приведены в таблице \ref{app_set:urovnn}.

\begin{figure}[H]
\centering
\includegraphics[width=1\textwidth,height=1\textheight,keepaspectratio]{\fbpath/06.02 ДЗАТ/ПО УРОВ НН/img/UROV_NN.pdf}
\captionof{figure}{Функциональная схема логики ПО УРОВ НН}\label{pic:UROVNN}
\end{figure}



\end{enumerate}